\documentclass[12pt,letterpaper,oneside]{report}

% Kompilator w Overleaf: pdfLaTeX
\usepackage[utf8]{inputenc}
\usepackage[T1]{fontenc}
\usepackage[polish]{babel}
\usepackage{polski}

\usepackage[letterpaper,left=2.0cm,right=2.0cm,top=2.0cm,bottom=2.0cm]{geometry}
\usepackage{microtype}
\setlength{\parindent}{1.25em}
\setlength{\parskip}{0pt}
\setlength{\emergencystretch}{2em}

\usepackage{amsmath,amssymb,mathtools,bm}
\usepackage{graphicx}
\graphicspath{{figures/}}
\usepackage{booktabs,longtable,array,calc}
\usepackage{etoolbox}
\usepackage{caption}
\captionsetup{font=small,labelfont=normalfont,justification=justified,singlelinecheck=false}
\usepackage{float}
\usepackage{enumitem}
\setlist[itemize]{topsep=0.35em,itemsep=0.15em,leftmargin=1.6em}
\setlist[enumerate]{topsep=0.35em,itemsep=0.2em,leftmargin=1.9em}

\usepackage{titlesec}
\titleformat{\chapter}[display]
  {\normalfont\Large\bfseries}
  {\chaptername\ \thechapter}
  {0.85em}
  {}
\titlespacing*{\chapter}{0pt}{0pt}{2.2em}
\titleformat{\section}
  {\normalfont\large\bfseries}
  {\thesection}
  {0.9em}
  {}
\titlespacing*{\section}{0pt}{1.7em}{0.65em}
\titleformat{\subsection}
  {\normalfont\normalsize\bfseries}
  {\thesubsection}
  {0.8em}
  {}
\titlespacing*{\subsection}{0pt}{1.3em}{0.45em}

% W pracy wzorcowej numeracja podrozdziałów zaczyna się od 1 w każdym rozdziale.
\renewcommand{\thesection}{\arabic{section}}
\renewcommand{\thesubsection}{\thesection.\arabic{subsection}}
\numberwithin{equation}{chapter}
\numberwithin{figure}{chapter}
\numberwithin{table}{chapter}

\usepackage{tocloft}
\renewcommand{\contentsname}{Spis treści}
\renewcommand{\cfttoctitlefont}{\hfill\Large\bfseries}
\renewcommand{\cftaftertoctitle}{\hfill}
\setlength{\cftbeforetoctitleskip}{0.5cm}
\setlength{\cftaftertoctitleskip}{1.0cm}
\setlength{\cftbeforechapskip}{0.25em}
\setlength{\cftbeforesecskip}{0pt}

% Dane strony tytułowej - wystarczy zmienić poniższe linie.
\newcommand{\UniversityName}{Uniwersytet Jagielloński w Krakowie}
\newcommand{\FacultyName}{Wydział Fizyki, Astronomii i Informatyki Stosowanej}
\newcommand{\AuthorName}{Karina Nepravskaya}
\newcommand{\AlbumNumber}{1209730}
\newcommand{\SupervisorName}{tytuł, imię i nazwisko promotora}
\newcommand{\DepartmentName}{nazwa zakładu lub instytutu}
\newcommand{\YearOfSubmission}{2026}

\usepackage{hyperref}
\hypersetup{
  colorlinks=true,
  linkcolor=blue,
  citecolor=blue,
  urlcolor=blue,
  pdftitle={Analiza korelacji pędów elektronów w potrójnej jonizacji neonu},
  pdfauthor={\AuthorName}
}

% Skróty używane w tekście.
\newcommand{\Pe}{\ensuremath{P_e}}
\newcommand{\Neff}{\ensuremath{N_{\mathrm{eff}}}}
\newcommand{\xmax}{\ensuremath{x_{\max}}}
\newcommand{\direct}{\textit{direct}}
\newcommand{\delayed}{\textit{delayed}}

% Zwykły, pozbawiony ramki akapit wyróżniony - zgodny z oszczędnym stylem pracy wzorcowej.
\newenvironment{importantbox}[1]
  {\par\medskip\noindent\textbf{#1.}\quad}
  {\par\medskip}

% Pomocnicze definicje wymagane przez tabele wygenerowane przez Pandoc.
\providecommand{\tightlist}{\setlength{\itemsep}{0pt}\setlength{\parskip}{0pt}}
\makeatletter
\patchcmd\longtable{\par}{\if@noskipsec\mbox{}\fi\par}{}{}
\makeatother

\pagestyle{plain}
\setcounter{secnumdepth}{2}
\setcounter{tocdepth}{2}

\begin{document}

\begin{titlepage}
\thispagestyle{empty}

\noindent
{\bfseries \UniversityName}\par
\noindent
\FacultyName

\vspace{2.7cm}

\begin{center}
\centering
\LARGE Karina Nepravskaya

\vspace{0.35cm}
{\large Nr albumu: \AlbumNumber}
\end{center}

\vspace{2.0cm}

\begin{center}
{\LARGE\bfseries
Analiza korelacji pędów elektronów\\[0.2em]
w potrójnej jonizacji neonu\\[0.2em]
w silnym polu laserowym\par}

\vspace{1.2cm}

{\normalsize Notatka metodologiczna}
\end{center}

\vfill

\begin{flushright}
\hfill
\begin{minipage}{0.48\textwidth}
\raggedleft
Praca wykonana pod kierunkiem\\
\SupervisorName\\
\DepartmentName
\end{minipage}
\end{flushright}

\vfill

\begin{center}
Kraków \YearOfSubmission
\end{center}
\end{titlepage}

\pagenumbering{arabic}
\setcounter{page}{2}
\tableofcontents
\clearpage
\phantomsection
\addcontentsline{toc}{chapter}{Streszczenie}
\begin{center}
  {\bfseries Streszczenie}
\end{center}
\vspace{1.2em}

Celem analizy jest wyodrębnienie i ilościowe opisanie korelacji między
końcowymi składowymi pędu trzech elektronów emitowanych w potrójnej
jonizacji atomu neonu. Punktem wyjścia była analiza głównych składowych
przeprowadzona w dziewięciowymiarowej przestrzeni składowych pędu. PCA
wskazała, że dynamikę wzdłuż osi polaryzacji lasera można rozdzielić na
jeden kierunek kolektywny, odpowiadający zgodnemu ruchowi wszystkich
elektronów, oraz dwa niezależne kierunki względne, opisujące
nierównomierność podziału pędu. Zamiast posługiwać się bezpośrednio
zależnymi od zbioru danymi wektorami własnymi PCA, wprowadzono stałą,
ortonormalną bazę Jacobiego--Helmerta. Umożliwia ona porównywanie
kanałów direct i delayed w tym samym układzie współrzędnych oraz
zachowuje pełną informację o trzech podłużnych składowych pędu.

Współrzędne liniowe uzupełniono obserwablami zaczerpniętymi z
reprezentacji simpleksowej: udziałami bezwzględnego pędu, efektywną
liczbą elektronów uczestniczących w podziale pędu, największym udziałem
jednego elektronu, entropią podziału, sektorami znaków oraz
współrzędnymi logarytmicznych ilorazów. Taki zestaw oddziela trzy różne
aspekty zdarzenia: skalę i kierunek ruchu kolektywnego, stopień
nierównowagi względnej oraz strukturę podziału wartości bezwzględnych i
znaków pędów. Dla danych przy natężeniu \ensuremath{1{,}7\,\mathrm{PW/cm^2}} kanał direct
charakteryzuje się silniejszymi dodatnimi korelacjami parami, mniejszą
wartością względnej nierównowagi K i bardziej równomiernym podziałem
pędu niż kanał delayed. Wyniki te są zgodne z obrazem bardziej
kolektywnej, niemal jednoczesnej emisji w kanale direct oraz bardziej
rozdzielonej czasowo dynamiki w kanale delayed, przy czym same pędy
końcowe nie stanowią samodzielnego dowodu mechanizmu czasowego.

\begin{importantbox}{Najważniejsza idea metody}
PCA służy tutaj jako narzędzie rozpoznania struktury danych, a nie jako
ostateczny układ współrzędnych. Ostateczna baza jest ustalona przez
kinematykę układu wielu cząstek. Dzięki temu osie nie obracają się przy
zmianie kanału jonizacji, natężenia pola ani liczebności próbki.
\end{importantbox}


\clearpage
\thispagestyle{empty}
\null
\clearpage
\chapter{Wstęp}

Wielokrotna jonizacja w silnym polu laserowym może mieć charakter skorelowany: końcowy ruch elektronów zależy nie tylko od pola, lecz także od energii przekazywanej pomiędzy nimi. W przypadku potrójnej jonizacji obserwujemy jednocześnie trzy wektory pędu. Pojawia się więc problem, jak rozdzielić wspólny ruch elektronów od różnic wewnątrz trójki i jak porównywać kilka klas zdarzeń bez sprowadzania wielowymiarowego rozkładu do jednej niejednoznacznej projekcji.

Szczególnie ważne są składowe wzdłuż osi polaryzacji. Emisja elektronów w zbliżonej fazie pola może prowadzić do podobnych wartości i znaków pędu, natomiast emisja rozdzielona w czasie może zwiększać różnice. Rozkład jednego elektronu ani sama suma pędów nie wystarczają do opisania obu sytuacji. Przykładowo konfiguracje $(2,2,2)$ i $(5,1,0)$ mają tę samą sumę, lecz zupełnie inny podział.

Dla trzech elektronów podział wartości bezwzględnych można przedstawić na trójkącie simpleksowym. Dla czterech elektronów powstaje jednak tetraedr, a dla $N$ elektronów simpleks o wymiarze $N-1$. Wraz ze wzrostem liczby cząstek bezpośrednia wizualizacja przestaje być czytelna. Potrzebny jest zatem zestaw wielkości, które opisują oddzielnie ruch kolektywny, ruch względny, równomierność podziału i kierunek emisji.

Punktem wyjścia pracy jest Analiza Głównych Składowych (PCA) zastosowana do dziewięciu elektronowych składowych pędu. PCA służy do rozpoznania struktury danych, ale jej osie zależą od konkretnej macierzy kowariancji. Wynik wykorzystano więc jako motywację do wprowadzenia jednej stałej bazy kolektywnej i względnej, wspólnej dla wszystkich analizowanych zbiorów.

Celem pracy jest opracowanie spójnego sposobu analizy korelacji końcowych pędów elektronów w potrójnej jonizacji neonu oraz porównanie kanałów \direct{} i \delayed{} dla natężeń $1{,}0$, $1{,}3$, $1{,}6$ i $1{,}7\,\mathrm{PW/cm^2}$. Sprawdzono, czy PCA wyróżnia kierunek kolektywny, czy różnice pomiędzy kanałami utrzymują się przy zmianie natężenia oraz które elementy opisu można uogólnić na większą liczbę elektronów.

Rozdział drugi przedstawia fizyczny kontekst jonizacji wielokrotnej. W rozdziale trzecim opisano model i dane. Rozdział czwarty zawiera metody analizy, piąty wyniki dla wszystkich natężeń, a szósty ich interpretację i ograniczenia. Ostatni rozdział podsumowuje najważniejsze wnioski.


\chapter{Jonizacja wielokrotna w silnym polu laserowym}

\section{Tunelowanie i recollision}

Silne pole elektryczne odkształca potencjał atomowy i umożliwia elektronowi przejście przez barierę do kontinuum. Uwolniony elektron jest następnie przyspieszany przez oscylujące pole. Po zmianie jego zwrotu może powrócić w pobliże jonu i przekazać energię pozostałym elektronom. Takie zderzenie powrotne nazywa się \textit{recollision}.

Jeżeli elektrony są usuwane kolejno głównie przez pole, jonizacja ma charakter sekwencyjny. Jeżeli emisję kilku elektronów uruchamia transfer energii podczas recollision, proces jest niesekwencyjny. Końcowe pędy zawierają wtedy informację o dynamice wielu cząstek, ponieważ zależą zarówno od wspólnego zderzenia, jak i późniejszego przyspieszania w polu.

Potrójna jonizacja jest trudniejsza od podwójnej: energia zostaje rozdzielona pomiędzy trzy elektrony, które nie muszą opuścić atomu jednocześnie. Pełne obliczenia kwantowe w trzech wymiarach są kosztowne, dlatego do analizy takich procesów stosuje się modele półklasyczne pozwalające śledzić czasy jonizacji i pojedyncze trajektorie.

\section{Kanały direct i delayed}

W kanale \direct{} wszystkie trzy elektrony jonizują krótko po głównym zderzeniu powrotnym. W kanale \delayed{} recollision prowadzi do szybkiej emisji dwóch elektronów, natomiast trzeci opuszcza atom z opóźnieniem. Klasyfikacja wynika z historii czasowej trajektorii, a nie z pędów końcowych.

Kanały są klasami statystycznymi, a nie ścisłymi wzorcami pędu. Można jednak oczekiwać, że elektrony emitowane w podobnej fazie pola częściej będą miały zgodne znaki i mniejsze różnice pędów. Elektron jonizujący później może zostać przyspieszony w innej fazie, zwiększając nierównowagę wewnątrz zdarzenia.

\section{Korelacje pędowe}

Suma pędów opisuje ruch całej trójki, ale nie jej wewnętrzny podział. Udziały wartości bezwzględnych opisują proporcje, lecz usuwają znaki. Korelacje parami wykrywają zależność liniową, ale nie stanowią pełnego opisu trójelektronowego. Z tego powodu żadna pojedyncza obserwabla nie jest wystarczająca. W pracy wykorzystano kilka uzupełniających się reprezentacji, z których każda odpowiada na inne pytanie.


\chapter{Model i dane}

\section{Model ECBB}

Dane pochodzą z trójwymiarowego modelu półklasycznego ECBB. Propagowane są równocześnie jon i trzy aktywne elektrony. Jeden elektron rozpoczyna propagację po tunelowaniu, a dwa są początkowo związane. Model zachowuje pełne oddziaływanie kulombowskie dla par obejmujących elektron swobodny lub powracający. Dla pary elektronów związanych stosowany jest potencjał efektywny ograniczający sztuczną autojonizację klasycznego układu~\cite{Emmanouilidou2023}.

Model pozwala śledzić zderzenia i czasy jonizacji, ale nie jest pełnym opisem kwantowym. Nie obejmuje wszystkich efektów interferencji ani antysymetryzacji. Jego zaletą jest możliwość analizy dużej liczby pełnowymiarowych trajektorii i powiązania obserwabli końcowych z historią procesu.

\section{Parametry i struktura plików}

Dla natężeń $1{,}0$, $1{,}3$ i $1{,}6\,\mathrm{PW/cm^2}$ zastosowano impuls o długości fali $800\,\mathrm{nm}$ i czasie trwania natężenia $25\,\mathrm{fs}$. Pole jest spolaryzowane wzdłuż osi $z$. Analizę uzupełniono zbiorem przy $1{,}7\,\mathrm{PW/cm^2}$ zapisanym w tym samym formacie.

Każdy wiersz zawiera końcowe składowe pędu jonu i trzech elektronów:
\begin{equation}
(p_{1x},p_{1y},p_{1z},p_{2x},p_{2y},p_{2z},
p_{3x},p_{3y},p_{3z},p_{4x},p_{4y},p_{4z}).
\end{equation}
Cząstka 1 jest jonem, elektrony 2 i 3 są początkowo związane, a elektron 4 jest początkowo tunelujący. Numery opisują historię klasycznej trajektorii, a nie fundamentalne rozróżnienie identycznych elektronów.

\begin{table}[htbp]
\centering
\caption{Liczby zdarzeń w analizowanych plikach.}
\label{tab:counts-all}
\begin{tabular}{crr}
\toprule
$I\,[\mathrm{PW/cm^2}]$ & \direct & \delayed\\
\midrule
1,0 & 2042 & 2097\\
1,3 & 6198 & 5837\\
1,6 & 8799 & 7560\\
1,7 & 1695 & 1601\\
\bottomrule
\end{tabular}
\end{table}

PCA obejmuje dziewięć składowych elektronowych. Dalsza analiza skupia się na $p_2=p_{2z}$, $p_3=p_{3z}$ i $p_4=p_{4z}$. Każdy zbiór przetwarzany jest według tej samej procedury, a kanały są porównywane wyłącznie przy tym samym natężeniu.

\section{Ograniczenia danych}

Wyniki dotyczą symulacji dla określonych parametrów impulsu. Klasy \direct{} i \delayed{} są nadane na podstawie historii trajektorii. Pędy końcowe pokazują, jak klasyfikacja przejawia się w obserwablach asymptotycznych, ale nie służą do niezależnego wyznaczania etykiet.

Wielkości odwołujące się do elektronu 4 zachowują informację o historii modelu, lecz zależą od numeracji. Dlatego obok nich stosowane są miary symetryczne względem permutacji elektronów.


\chapter{Metody analizy}
\section{Analiza Głównych Składowych}

\section{Przygotowanie danych do PCA}

W notebooku 05\_PCA\_clean.ipynb PCA została przeprowadzona w przestrzeni
dziewięciu składowych pędu elektronów. Dla zbioru zawierającego $N$
zdarzeń dane zapisujemy w postaci macierzy
\begin{equation}
\mathbf X=
\begin{pmatrix}
(\mathbf x^{(1)})^{\mathsf T}\\
(\mathbf x^{(2)})^{\mathsf T}\\
\vdots\\
(\mathbf x^{(N)})^{\mathsf T}
\end{pmatrix}
\in \mathbb{R}^{N\times 9},
\label{eq:pca-data-matrix}
\end{equation}
gdzie $\mathbf x^{(i)}$ jest dziewięciowymiarowym wektorem składowych
pędu elektronów dla $i$-tego zdarzenia.

Przed wykonaniem PCA każda kolumna została wystandaryzowana. Dla
$j$-tej zmiennej stosujemy transformację
\begin{equation}
z_{ij}=\frac{x_{ij}-\overline{x}_j}{s_j},
\label{eq:pca-standardization}
\end{equation}
gdzie
\begin{equation}
\overline{x}_j=\frac{1}{N}\sum_{i=1}^{N}x_{ij},
\qquad
s_j^2=\frac{1}{N}\sum_{i=1}^{N}
\left(x_{ij}-\overline{x}_j\right)^2.
\label{eq:pca-mean-std}
\end{equation}
W wyniku otrzymujemy macierz $\mathbf Z$ o kolumnach mających średnią
równą zero i jednostkową skalę. W implementacji transformację tę
wykonuje klasa \texttt{StandardScaler}. Dzięki standaryzacji wynik PCA
nie jest zdominowany wyłącznie przez zmienne o największej bezwzględnej
wariancji.

PCA opiera się na diagonalizacji macierzy kowariancji danych
wystandaryzowanych
\begin{equation}
\mathbf C=\frac{1}{N-1}\mathbf Z^{\mathsf T}\mathbf Z.
\label{eq:pca-covariance}
\end{equation}
Ponieważ wejściowe zmienne zostały wystandaryzowane, macierz
$\mathbf C$ opisuje w tym przypadku przede wszystkim ich strukturę
korelacyjną. Kierunki głównych składowych są wektorami własnymi tej
macierzy i spełniają równanie
\begin{equation}
\mathbf C\mathbf v_k=\lambda_k\mathbf v_k,
\qquad
\lambda_1\geq\lambda_2\geq\ldots\geq\lambda_9,
\label{eq:pca-eigenproblem}
\end{equation}
przy czym wektory $\mathbf v_k$ są ortonormalne,
$\mathbf v_k^{\mathsf T}\mathbf v_l=\delta_{kl}$. Współczynniki
wektora $\mathbf v_k$ określają udział poszczególnych zmiennych w osi
$\mathrm{PC}k$. Wartości tych współczynników odnoszą się do przestrzeni
zmiennych wystandaryzowanych.

Dla pojedynczego zdarzenia współrzędna wzdłuż $k$-tej głównej składowej
jest rzutem wystandaryzowanego wektora na odpowiedni wektor własny,
\begin{equation}
t_{ik}=(\mathbf z^{(i)})^{\mathsf T}\mathbf v_k
      =\sum_{j=1}^{9}z_{ij}v_{jk}.
\label{eq:pca-score}
\end{equation}
Wartość własna $\lambda_k$ jest wariancją danych wzdłuż tej osi.
Udział wariancji wyjaśnianej przez daną składową definiujemy jako
\begin{equation}
r_k=\frac{\lambda_k}{\sum_{j=1}^{9}\lambda_j},
\label{eq:pca-explained-variance}
\end{equation}
a skumulowany udział pierwszych $m$ składowych jako
\begin{equation}
R_m=\sum_{k=1}^{m}r_k.
\label{eq:pca-cumulative-variance}
\end{equation}

Wynikiem PCA są zatem dwa rodzaje informacji wykorzystywane w dalszej
analizie: orientacja osi, określona przez współczynniki wektorów
$\mathbf v_k$, oraz odpowiadający każdej osi udział wyjaśnianej
wariancji $r_k$. W analizowanym problemie szczególne znaczenie mają te
osie, których współczynniki są skupione na składowych podłużnych
$p_{2z}$, $p_{3z}$ i $p_{4z}$.

PCA zwraca dziewięć ortogonalnych kierunków. W analizowanym zbiorze
direct pierwsza składowa wyjaśniała około 30,0\% całkowitej wariancji.
Składowe PC8 i PC9 wyjaśniały odpowiednio około 1,90\% i 1,46\%. Ich
mały udział nie oznacza, że są zbędne: są to kierunki o małej wariancji,
które domykają pełny trójwymiarowy podobszar ruchu wzdłuż osi z i
opisują różnice pomiędzy elektronami.

\section{Kierunki podłużne wskazane przez PCA}

Współczynniki PC1 dla składowych \ensuremath{p_{2z}}, \ensuremath{p_{3z}} i \ensuremath{p_{4z}} są prawie równe i mają
ten sam znak. Oznacza to, że PC1 odpowiada przede wszystkim wspólnemu
przesunięciu pędów wszystkich elektronów wzdłuż osi polaryzacji. Dla
zbioru direct otrzymano w przybliżeniu

\begin{equation}
\mathrm{PC1}\approx 0{,}580\,p_{2z}+0{,}577\,p_{3z}+0{,}572\,p_{4z}
\label{eq:4}
\end{equation}

Jest to numeryczny odpowiednik kierunku (1,1,1), czyli sumy pędów
elektronów. Składowe PC8 i PC9 mają natomiast współczynniki o
przeciwnych znakach, a ich suma jest bliska zeru. Opisują zatem
wewnętrzne odchylenia od ruchu kolektywnego. Dla kanałów direct i
delayed orientacja tych dwóch kierunków jest inna, mimo że w obu
przypadkach rozpinają one tę samą płaszczyznę względną.

\begin{longtable}[]{@{}
  >{\raggedright\arraybackslash}p{(\columnwidth - 8\tabcolsep) * \real{0.2000}}
  >{\raggedright\arraybackslash}p{(\columnwidth - 8\tabcolsep) * \real{0.2000}}
  >{\raggedright\arraybackslash}p{(\columnwidth - 8\tabcolsep) * \real{0.2000}}
  >{\raggedright\arraybackslash}p{(\columnwidth - 8\tabcolsep) * \real{0.2000}}
  >{\raggedright\arraybackslash}p{(\columnwidth - 8\tabcolsep) * \real{0.2000}}@{}}
\caption{Współczynniki składowych z w trzech kierunkach PCA zdominowanych przez ruch podłużny. Znak wektora własnego PCA jest umowny.}\label{tab:2}\\
\toprule\noalign{}
\begin{minipage}[b]{\linewidth}\raggedright
\textbf{Kanał}
\end{minipage} & \begin{minipage}[b]{\linewidth}\raggedright
\textbf{PC}
\end{minipage} & \begin{minipage}[b]{\linewidth}\raggedright
\textbf{w(\ensuremath{p_{2z}})}
\end{minipage} & \begin{minipage}[b]{\linewidth}\raggedright
\textbf{w(\ensuremath{p_{3z}})}
\end{minipage} & \begin{minipage}[b]{\linewidth}\raggedright
\textbf{w(\ensuremath{p_{4z}})}
\end{minipage} \\
\midrule\noalign{}
\endhead
\bottomrule\noalign{}
\endlastfoot
direct & PC1 & 0,580 & 0,577 & 0,572 \\
direct & PC8 & \ensuremath{-}0,283 & \ensuremath{-}0,517 & 0,808 \\
direct & PC9 & 0,763 & \ensuremath{-}0,631 & \ensuremath{-}0,136 \\
delayed & PC1 & 0,563 & 0,535 & 0,619 \\
delayed & PC8 & \ensuremath{-}0,610 & 0,764 & \ensuremath{-}0,105 \\
delayed & PC9 & \ensuremath{-}0,532 & \ensuremath{-}0,330 & 0,767 \\
\end{longtable}

W notebooku 05 jako pierwszą interpretowalną parametryzację
zaproponowano trzy wielkości: \ensuremath{P_{\mathrm{col}}} = \ensuremath{p_{2z}}+\ensuremath{p_{3z}}+\ensuremath{p_{4z}}, \ensuremath{A_3} = \ensuremath{p_{3z}}\ensuremath{-}(\ensuremath{p_{2z}}+\ensuremath{p_{4z}})/2
oraz \ensuremath{R_{24}} = \ensuremath{p_{2z}}\ensuremath{-}\ensuremath{p_{4z}}. Odpowiadające im wektory są wzajemnie ortogonalne,
choć nie są znormalizowane. Pierwszy z nich opisuje ruch wspólny, a dwa
pozostałe tworzą bazę płaszczyzny odchyleń względnych.

\section{Dlaczego nie używać bezpośrednio osi PCA}

Wektory PCA są wyznaczane osobno z macierzy kowariancji danego zbioru.
Ich orientacja zmienia się więc wraz z kanałem jonizacji, natężeniem
pola, kryteriami selekcji i statystyką próbki. Ponadto znak każdego
wektora własnego jest całkowicie umowny, a dwa kierunki o zbliżonych
wartościach własnych mogą obracać się w swojej wspólnej podprzestrzeni
bez istotnej zmiany jakości rekonstrukcji. Bezpośrednie porównywanie
współrzędnych PC8 i PC9 pomiędzy różnymi zbiorami mogłoby zatem
prowadzić do sztucznych różnic wynikających z obrotu osi, a nie z fizyki
procesu.

\begin{importantbox}{Wniosek z PCA}
PCA wiarygodnie identyfikuje strukturę „jeden kierunek kolektywny + dwa
kierunki względne''. Nie dostarcza jednak naturalnej, wspólnej dla
wszystkich zbiorów orientacji dwóch osi względnych. Tę orientację należy
ustalić na podstawie symetrii i kinematyki układu.
\end{importantbox}

\section{Stała baza kolektywna i względna}

\section{Definicja transformacji}

Dla każdego zdarzenia rozpatrzmy wektor trzech podłużnych składowych
pędu

\begin{equation}
\mathbf{p}_z=\bigl(p_2,p_3,p_4\bigr)^{\mathsf T}
\label{eq:5}
\end{equation}

Wprowadzamy ortonormalną bazę Helmerta, równoważną standardowym
współrzędnym Jacobiego dla trzech cząstek o jednakowych masach:

\begin{equation}
\mathbf e_0=\frac{(1,1,1)^{\mathsf T}}{\sqrt3},\qquad \mathbf e_1=\frac{(1,-1,0)^{\mathsf T}}{\sqrt2},\qquad \mathbf e_2=\frac{(1,1,-2)^{\mathsf T}}{\sqrt6}
\label{eq:6}
\end{equation}

Rzuty wektora \ensuremath{\mathbf p_z} na tę bazę definiują trzy nowe współrzędne:

\begin{equation}
Q=\frac{p_2+p_3+p_4}{\sqrt3}
\label{eq:7}
\end{equation}

\begin{equation}
\xi_1=\frac{p_2-p_3}{\sqrt2}
\label{eq:8}
\end{equation}

\begin{equation}
\xi_2=\frac{p_2+p_3-2p_4}{\sqrt6}
\label{eq:9}
\end{equation}

\begin{longtable}[]{@{}
  >{\raggedright\arraybackslash}p{(\columnwidth - 8\tabcolsep) * \real{0.2000}}
  >{\raggedright\arraybackslash}p{(\columnwidth - 8\tabcolsep) * \real{0.2000}}
  >{\raggedright\arraybackslash}p{(\columnwidth - 8\tabcolsep) * \real{0.2000}}
  >{\raggedright\arraybackslash}p{(\columnwidth - 8\tabcolsep) * \real{0.2000}}
  >{\raggedright\arraybackslash}p{(\columnwidth - 8\tabcolsep) * \real{0.2000}}@{}}
\caption{Macierz ortonormalnej transformacji Jacobiego--Helmerta.}\label{tab:3}\\
\toprule\noalign{}
\begin{minipage}[b]{\linewidth}\raggedright
\textbf{Współrzędna}
\end{minipage} & \begin{minipage}[b]{\linewidth}\raggedright
\textbf{\ensuremath{p_2}}
\end{minipage} & \begin{minipage}[b]{\linewidth}\raggedright
\textbf{\ensuremath{p_3}}
\end{minipage} & \begin{minipage}[b]{\linewidth}\raggedright
\textbf{\ensuremath{p_4}}
\end{minipage} & \begin{minipage}[b]{\linewidth}\raggedright
\textbf{Interpretacja}
\end{minipage} \\
\midrule\noalign{}
\endhead
\bottomrule\noalign{}
\endlastfoot
Q & \ensuremath{1/\sqrt{3}} & \ensuremath{1/\sqrt{3}} & \ensuremath{1/\sqrt{3}} & wspólny podłużny tryb pędu \\
\ensuremath{\xi_1} & \ensuremath{1/\sqrt{2}} & \ensuremath{-}\ensuremath{1/\sqrt{2}} & 0 & różnica między elektronami 2 i 3 \\
\ensuremath{\xi_2} & \ensuremath{1/\sqrt{6}} & \ensuremath{1/\sqrt{6}} & \ensuremath{-2/\sqrt{6}} & para (2,3) względem elektronu 4 \\
\end{longtable}

Współrzędna Q jest znormalizowaną sumą pędów. W notebooku do
wizualizacji użyto również wielkości \Pe{} = \ensuremath{p_2}+\ensuremath{p_3}+\ensuremath{p_4}, dla której
zachodzi \Pe{} = \ensuremath{\sqrt{3}} Q. Współrzędna \ensuremath{\xi_1} opisuje asymetrię w parze
elektronów początkowo związanych. Współrzędna \ensuremath{\xi_2} porównuje średni ruch
tej pary z ruchem elektronu oznaczonego jako początkowo tunelujący.

Baza zastosowana w notebooku 07 jest inną orientacją tej samej
płaszczyzny względnej co para \ensuremath{A_3} i \ensuremath{R_{24}} z notebooka 05. Oba warianty są
matematycznie równoważne: przejście pomiędzy nimi jest obrotem w
dwuwymiarowej przestrzeni prostopadłej do kierunku (1,1,1). Wybrana
orientacja ma jednak czytelniejszą interpretację związaną z rolami
początkowymi elektronów.

\section{Moduł i kąt nierównowagi względnej}

Dwie współrzędne względne można traktować jako kartezjańskie osie
płaszczyzny podziału pędu. Ich promień definiuje miarę całkowitej
nierównowagi:

\begin{equation}
K=\sqrt{\xi_1^2+\xi_2^2}
\label{eq:10}
\end{equation}

Wielkość K ma wymiar pędu. Przyjmuje wartość zero wtedy i tylko wtedy,
gdy \ensuremath{p_2}=\ensuremath{p_3}=\ensuremath{p_4}. Można ją zapisać w dwóch równoważnych postaciach:

\begin{equation}
K^2=\sum_{i=2}^{4}(p_i-\bar p)^2=\frac{(p_2-p_3)^2+(p_2-p_4)^2+(p_3-p_4)^2}{3}
\label{eq:11}
\end{equation}

\begin{equation}
\bar p=\frac{p_2+p_3+p_4}{3}
\label{eq:12}
\end{equation}

Z równania~\eqref{eq:11} wynika, że K jest niezmiennicze na dowolną permutację
elektronów. Jest to ważne w kontekście nierozróżnialności elektronów:
chociaż same osie \ensuremath{\xi_1} i \ensuremath{\xi_2} zależą od przyjętej etykiety, odległość od
początku układu względnego nie zależy od numeracji.

Notebook oblicza również kąt \ensuremath{\varphi} = atan2(\ensuremath{\xi_2},\ensuremath{\xi_1}). Wielkość K mówi, jak duża
jest nierównowaga, natomiast \ensuremath{\varphi} wskazuje, w którym kierunku płaszczyzny
względnej leży zdarzenie, czyli który kontrast pędów dominuje.
Zastąpienie pary (\ensuremath{\xi_1},\ensuremath{\xi_2}) samym K usuwa informację kątową; dlatego pełny
histogram dwuwymiarowy i rozkład K należy traktować jako obserwable
komplementarne.

\section{Zachowanie normy i odwracalność}

Ponieważ macierz Helmerta jest ortogonalna, transformacja zachowuje
normę euklidesową:

\begin{equation}
p_2^2+p_3^2+p_4^2=Q^2+\xi_1^2+\xi_2^2
\label{eq:13}
\end{equation}

Nie jest to redukcja wymiaru, lecz zmiana bazy. Trzy pierwotne składowe
można dokładnie odtworzyć ze współrzędnych nowych:

\begin{equation}
p_2=\frac{Q}{\sqrt3}+\frac{\xi_1}{\sqrt2}+\frac{\xi_2}{\sqrt6}
\label{eq:14}
\end{equation}

\begin{equation}
p_3=\frac{Q}{\sqrt3}-\frac{\xi_1}{\sqrt2}+\frac{\xi_2}{\sqrt6}
\label{eq:15}
\end{equation}

\begin{equation}
p_4=\frac{Q}{\sqrt3}-\frac{2\xi_2}{\sqrt6}
\label{eq:16}
\end{equation}

\section{Obserwable uzupełniające}

\subsection{Odchylenia od średniego pędu i bilans z jonem}

Dla każdego zdarzenia obliczono odchylenia \ensuremath{q_2}=\ensuremath{p_2}\ensuremath{-}\ensuremath{\bar p}, \ensuremath{q_3}=\ensuremath{p_3}\ensuremath{-}\ensuremath{\bar p} i \ensuremath{q_4}=\ensuremath{p_4}\ensuremath{-}\ensuremath{\bar p}. Z
definicji zachodzi \ensuremath{q_2}+\ensuremath{q_3}+\ensuremath{q_4}=0. Wielkości te są bezpośrednią
reprezentacją ruchu wewnętrznego po odjęciu składowej kolektywnej.
Dodatkowo wyznaczono sumę \ensuremath{p_{1z}}+\ensuremath{p_{2z}}+\ensuremath{p_{3z}}+\ensuremath{p_{4z}}, która może służyć jako
kontrola zachowania pędu całego układu wzdłuż osi z i jako diagnostyka
wpływu procedury numerycznej lub selekcji zdarzeń.

\subsection{Udziały bezwzględnego pędu}

Aby oddzielić sam sposób podziału od ogólnej skali pędu, zdefiniowano
normę \ensuremath{L^1}

\begin{equation}
A=|p_2|+|p_3|+|p_4|
\label{eq:17}
\end{equation}

oraz udziały \ensuremath{x_2}, \ensuremath{x_3} i \ensuremath{x_4} zgodnie z równaniem~\eqref{eq:2}. Na ich podstawie
obliczono trzy skalarne miary:

\begin{itemize}
\item
  \Neff{} = 1/(\ensuremath{x_2}\textsuperscript{2}+\ensuremath{x_3}\textsuperscript{2}+\ensuremath{x_4}\textsuperscript{2}) --- efektywna liczba elektronów
  uczestniczących w podziale bezwzględnego pędu. Dla trzech elektronów 1
  \ensuremath{\leq} \Neff{} \ensuremath{\leq} 3. Wartość 1 oznacza dominację jednego elektronu, a wartość
  3 idealnie równy podział.
\item
  \xmax{} = max(\ensuremath{x_2},\ensuremath{x_3},\ensuremath{x_4}) --- największy udział pojedynczego elektronu.
  Zakres wynosi od 1/3 do 1; mniejsze wartości odpowiadają bardziej
  równomiernemu podziałowi.
\item
  H = \ensuremath{-}\ensuremath{\sum_i}\ensuremath{x_i} ln(\ensuremath{x_i})/ln 3 --- znormalizowana entropia udziałów. Wartość 0
  odpowiada pełnej koncentracji pędu w jednym elektronie, a wartość 1
  równemu podziałowi.
\end{itemize}

Miary \Neff{}, \xmax{} i H są bezwymiarowe i niezmiennicze na permutacje
elektronów. Uzupełniają K, które ma wymiar pędu i rośnie zarówno wskutek
silnej nierównowagi, jak i zwiększenia ogólnej skali pędów względnych.

\subsection{Sektory znaków}

Normalizacja przez wartości bezwzględne usuwa informację o kierunku
emisji. Z tego powodu każdemu zdarzeniu przypisano sektor znaków
określony przez znaki trójki (\ensuremath{p_2},\ensuremath{p_3},\ensuremath{p_4}). Dla trzech elektronów występuje
osiem sektorów: \ensuremath{+++}, \ensuremath{++-}, \ensuremath{+-+}, \ensuremath{-++}, \ensuremath{+--}, \ensuremath{-+-}, \ensuremath{--+} i \ensuremath{---}. Sektory \ensuremath{+++} i
\ensuremath{---} oznaczają zgodny kierunek wszystkich elektronów, natomiast pozostałe
sektory odpowiadają emisji mieszanej.

Pełna etykieta sektora zachowuje informację o rolach początkowych
elektronów, ale zależy od numeracji. W analizach nastawionych wyłącznie
na obserwable permutacyjnie symetryczne można dodatkowo grupować
zdarzenia według liczby dodatnich pędów: 3, 2, 1 lub 0.

\subsection{Współrzędne ilr dla danych kompozycyjnych}

Udziały \ensuremath{x_i} tworzą dane kompozycyjne, dla których naturalnymi
współrzędnymi są logarytmy ilorazów. W notebooku zastosowano
transformację isometric log-ratio:

\begin{equation}
y_1=\frac{1}{\sqrt2}\ln\!\left(\frac{x_2}{x_3}\right)
\label{eq:18}
\end{equation}

\begin{equation}
y_2=\frac{1}{\sqrt6}\ln\!\left(\frac{x_2x_3}{x_4^2}\right)
\label{eq:19}
\end{equation}

Punkt (0,0) odpowiada równym udziałom \ensuremath{x_2}=\ensuremath{x_3}=\ensuremath{x_4}=1/3. Oddalanie się od
początku oznacza coraz większą nierównowagę względnych udziałów.
Transformacja ilr przedstawia wnętrze trójkąta simpleksowego na zwykłej
płaszczyźnie euklidesowej. Ponieważ logarytm nie jest określony dla
zera, w implementacji zastosowano małą dodatnią wartość regularizującą
przed obliczeniem ilorazów.

\subsection{Korelacje i moment trzeciego rzędu}

Dla każdej pary elektronów obliczono współczynnik korelacji Pearsona
pomiędzy pędami podłużnymi. Wysoka dodatnia korelacja oznacza, że dwa
elektrony zwykle zwiększają lub zmniejszają swój pęd w tym samym
kierunku. Współczynnik Pearsona opisuje jednak wyłącznie zależność
liniową i nie rozróżnia wszystkich struktur wielociałowych.

Jako pomocniczą miarę trójcząstkową wyznaczono wycentrowany moment
trzeciego rzędu

\begin{equation}
\kappa_{234}=\left\langle (p_2-\langle p_2\rangle)(p_3-\langle p_3\rangle)(p_4-\langle p_4\rangle)\right\rangle
\label{eq:20}
\end{equation}

Dla trzech różnych zmiennych jest on równy mieszanemu kumulantowi
trzeciego rzędu. Wielkość ta ma wymiar pędu do trzeciej potęgi, jest
wrażliwa na asymetrię rozkładu i skalę pędów, dlatego nie powinna być
interpretowana bez normalizacji lub oszacowania niepewności
statystycznej.

\section{Przebieg analizy i kontrola poprawności}
\section{Przebieg obliczeń w notebooku 07\_analiza\_3D}

Analiza została zorganizowana jako powtarzalny potok obliczeniowy. Dla
aktualnie aktywnej konfiguracji wykorzystano dwa zbiory przy natężeniu
I=\ensuremath{1{,}7\,\mathrm{PW/cm^2}}: direct oraz delayed. Pliki są odczytywane zarówno przy
separatorze przecinkowym, jak i białych znakach, a liczba kolumn jest
kontrolowana przed dalszym przetwarzaniem.

\begin{enumerate}
\def\labelenumi{\arabic{enumi}.}
\item
  Zdefiniowanie listy zbiorów danych wraz z natężeniem pola, kanałem
  jonizacji i ścieżką do pliku.
\item
  Odczyt dwunastu kolumn pędu, usunięcie pustych wierszy i wymuszenie
  typu zmiennoprzecinkowego.
\item
  Wydzielenie składowych \ensuremath{p_{2z}}, \ensuremath{p_{3z}} i \ensuremath{p_{4z}} oraz obliczenie \Pe{}, Q, \ensuremath{\xi_1}, \ensuremath{\xi_2},
  K i \ensuremath{\varphi}.
\item
  Obliczenie odchyleń \ensuremath{q_i}, sum kontrolnych oraz różnicy pomiędzy normą w
  bazie pierwotnej i bazie Jacobiego--Helmerta.
\item
  Wyznaczenie udziałów \ensuremath{x_i}, miar \Neff{}, \xmax{} i H, sektorów znaków oraz
  współrzędnych ilr.
\item
  Obliczenie różnic parami \ensuremath{d_{23}}, \ensuremath{d_{24}} i \ensuremath{d_{34}} oraz macierzy korelacji pędów
  podłużnych.
\item
  Połączenie wszystkich zbiorów w jedną tabelę combined, przy zachowaniu
  kolumn identyfikujących natężenie i kanał.
\item
  Budowa tabeli zbiorczej zawierającej liczebności, średnie, odchylenia
  standardowe, mediany, korelacje i moment trzeciego rzędu.
\item
  Wykonanie histogramów dwuwymiarowych i jednowymiarowych. W aktualnym
  notebooku użyto 120 przedziałów dla histogramów 2D, 80 przedziałów dla
  histogramów 1D oraz rozdzielczości zapisu 170 dpi.
\item
  Porównanie direct--delayed przy tym samym natężeniu oraz eksport
  danych przetworzonych, tabel podsumowujących, błędów rekonstrukcji i
  rysunków.
\end{enumerate}

\section{Kontrola poprawności transformacji}

Poprawność zmiany bazy została sprawdzona na kilku niezależnych
poziomach. Maksymalna różnica pomiędzy \ensuremath{p_2^2}+\ensuremath{p_3^2}+\ensuremath{p_4^2} a \ensuremath{Q^2}+\ensuremath{\xi_1^2}+\ensuremath{\xi_2^2} wyniosła
około \ensuremath{2{,}1\times 10^{-14}}. Maksymalna wartość \textbar \ensuremath{q_2}+\ensuremath{q_3}+\ensuremath{q_4}\textbar{} była
rzędu \ensuremath{1{,}8\times 10^{-15}}. Są to wartości odpowiadające precyzji obliczeń
zmiennoprzecinkowych.

Dodatkowo wykonano rekonstrukcję pędów na podstawie równań~\eqref{eq:14}--\eqref{eq:16}.
Maksymalne błędy bezwzględne wyniosły \ensuremath{1{,}78\times 10^{-15}} dla \ensuremath{p_2}, \ensuremath{1{,}55\times 10^{-15}} dla
\ensuremath{p_3} oraz \ensuremath{1{,}78\times 10^{-15}} dla \ensuremath{p_4}. Oznacza to, że współrzędne Q, \ensuremath{\xi_1} i \ensuremath{\xi_2}
zachowują pełną informację o trzech podłużnych składowych pędu. Utrata
informacji pojawia się dopiero wtedy, gdy para (\ensuremath{\xi_1},\ensuremath{\xi_2}) zostaje
zastąpiona samym promieniem K lub gdy udziały \ensuremath{x_i} są analizowane bez
sektorów znaków i bez skali A.

\begin{longtable}[]{@{}
  >{\raggedright\arraybackslash}p{(\columnwidth - 2\tabcolsep) * \real{0.5000}}
  >{\raggedright\arraybackslash}p{(\columnwidth - 2\tabcolsep) * \real{0.5000}}@{}}
\caption{Numeryczne testy ortogonalności i odwracalności transformacji.}\label{tab:4}\\
\toprule\noalign{}
\begin{minipage}[b]{\linewidth}\raggedright
\textbf{Test}
\end{minipage} & \begin{minipage}[b]{\linewidth}\raggedright
\textbf{Wartość maksymalna}
\end{minipage} \\
\midrule\noalign{}
\endhead
\bottomrule\noalign{}
\endlastfoot
\textbar \ensuremath{p_2^2}+\ensuremath{p_3^2}+\ensuremath{p_4^2}\ensuremath{-}\ensuremath{Q^2}\ensuremath{-}\ensuremath{\xi_1^2}\ensuremath{-}\ensuremath{\xi_2^2}\textbar{} & \ensuremath{\approx} \ensuremath{2{,}1\times 10^{-14}} \\
\textbar \ensuremath{q_2}+\ensuremath{q_3}+\ensuremath{q_4}\textbar{} & \ensuremath{\approx} \ensuremath{1{,}8\times 10^{-15}} \\
błąd rekonstrukcji \ensuremath{p_2} & \ensuremath{1{,}78\times 10^{-15}} \\
błąd rekonstrukcji \ensuremath{p_3} & \ensuremath{1{,}55\times 10^{-15}} \\
błąd rekonstrukcji \ensuremath{p_4} & \ensuremath{1{,}78\times 10^{-15}} \\
\end{longtable}

\section{Interpretacja zastosowanych wizualizacji}

\subsection{Płaszczyzna pędu względnego (\ensuremath{\xi_1}, \ensuremath{\xi_2})}

Histogram dwuwymiarowy na płaszczyźnie (\ensuremath{\xi_1},\ensuremath{\xi_2}) przedstawia pełną
strukturę względną po usunięciu ruchu kolektywnego. Początek układu
odpowiada równym pędom wszystkich elektronów. Odległość od początku jest
równa K, a kierunek określa typ asymetrii. Rozkład delayed jest wyraźnie
szerszy w obu kierunkach niż rozkład direct. Kanał direct tworzy
bardziej zwartą chmurę skupioną bliżej początku, co zapowiada mniejszą
średnią nierównowagę pędu.

\begin{figure}[htbp]
\centering
\includegraphics[width=0.96\textwidth]{figures/relative_plane.png}
\caption{Gęstość zdarzeń na płaszczyźnie względnej (\ensuremath{\xi_1},\ensuremath{\xi_2}): kanał delayed po lewej i direct po prawej. Osie są wspólne fizycznie, a nie wyznaczane osobno przez PCA.}
\label{fig:relative-plane}
\end{figure}

\subsection{Całkowity pęd \Pe{} a nierównowaga K}

Wykres (\Pe{},K) oddziela ruch środka pędu elektronów od wewnętrznego
podziału. W obu kanałach widoczna jest struktura po obu stronach \Pe{}=0,
związana z emisją w przeciwnych półcyklach pola. W kanale direct gęstość
skupia się przy większych wartościach \ensuremath{|P_e|} i
stosunkowo małym K. W kanale delayed rozkład K jest szerszy, a obszar
pośrednich wartości \Pe{} jest silniej wypełniony. Oznacza to, że duży
wspólny pęd i duża nierównowaga nie są tą samą cechą zdarzenia.

\begin{figure}[htbp]
\centering
\includegraphics[width=0.96\textwidth]{figures/P_vs_K.png}
\caption{Zależność pomiędzy całkowitym podłużnym pędem elektronów \Pe{} a modułem nierównowagi względnej K.}
\label{fig:p-vs-k}
\end{figure}

\subsection{Efektywna liczba elektronów \Neff{}}

Rozkład \Neff{} w kanale direct silnie narasta w pobliżu wartości 3,
czyli konfiguracji zbliżonych do równomiernego podziału bezwzględnego
pędu. Dla kanału delayed rozkład jest szerszy i zawiera znacznie większy
udział zdarzeń z \Neff{}\ensuremath{\approx}2 lub mniejszym. Jest to skalarne odtworzenie
informacji, która w trójkącie simpleksowym odpowiadałaby koncentracji
punktów odpowiednio w pobliżu środka albo w kierunku krawędzi i
wierzchołków.

\begin{figure}[htbp]
\centering
\includegraphics[width=0.96\textwidth]{figures/Neff.png}
\caption{Histogram efektywnej liczby elektronów uczestniczących w podziale bezwzględnego pędu.}
\label{fig:neff}
\end{figure}

\subsection{Sektory znaków}

Kanał direct jest zdominowany przez sektory \ensuremath{+++} i \ensuremath{---}, a sektory
mieszane mają niewielką liczebność. Oznacza to, że w większości zdarzeń
wszystkie trzy elektrony uzyskują końcowy pęd w tym samym kierunku
wzdłuż osi polaryzacji. W kanale delayed sektory zgodne nadal należą do
najliczniejszych, lecz zdarzenia o mieszanych znakach stanowią istotną
część rozkładu. Sama analiza wartości bezwzględnych nie ujawniłaby tej
różnicy.

\begin{figure}[htbp]
\centering
\includegraphics[width=0.96\textwidth]{figures/sign_sectors.png}
\caption{Liczebności ośmiu sektorów znaków (\ensuremath{p_{2z}},\ensuremath{p_{3z}},\ensuremath{p_{4z}}).}
\label{fig:sign-sectors}
\end{figure}

\subsection{Gęstość ilr udziałów bezwzględnych}

Współrzędne ilr przedstawiają ten sam znormalizowany podział, który w
reprezentacji simpleksowej znajduje się w trójkącie, lecz w zwykłym
układzie kartezjańskim. Kanał direct jest silnie skupiony w pobliżu
(0,0), co odpowiada udziałom zbliżonym do 1/3. Rozkład delayed jest
wyraźnie szerszy i bardziej rozciągnięty. Obraz jest zgodny z
histogramem \Neff{}, ale zachowuje również informację o kierunku
asymetrii udziałów.

\begin{figure}[htbp]
\centering
\includegraphics[width=0.96\textwidth]{figures/ilr.png}
\caption{Gęstość w przestrzeni ilr znormalizowanych udziałów \textbar \ensuremath{p_i}\textbar.}
\label{fig:ilr}
\end{figure}

\subsection{Macierze korelacji parami}

Macierz korelacji daje bezpośrednią liczbową ocenę zgodności zmian pędów
poszczególnych elektronów. W kanale direct wszystkie współczynniki poza
diagonalą przekraczają 0,81, co wskazuje na bardzo silną dodatnią
zależność liniową. W kanale delayed korelacja pomiędzy elektronami 2 i 3
wynosi jedynie około 0,21, a korelacje z elektronem 4 są umiarkowane,
około 0,44--0,46. Wynik ten jest spójny z większą szerokością rozkładu
względnego i częstszą emisją do mieszanych sektorów znaków.

\begin{figure}[htbp]
\centering
\includegraphics[width=0.96\textwidth]{figures/correlations.png}
\caption{Macierze korelacji Pearsona pomiędzy podłużnymi pędami elektronów.}
\label{fig:correlations}
\end{figure}

\chapter{Wyniki}

\section{Porównanie czterech natężeń}

Dla każdego natężenia kanał \direct{} ma szerszy rozkład całkowitego pędu $P_e$, lecz mniejszą średnią skalę ruchu względnego $K$. Wyniki trzech miar równomierności są zgodne: w \direct{} wartość bezwzględna pędu jest dzielona bardziej równomiernie.

\begin{table}[htbp]
\centering
\caption{Wyniki dla wszystkich analizowanych zbiorów.}
\label{tab:summary-all}
\small
\begin{tabular}{clrrrrr}
\toprule
$I$ & kanał & $\sigma(P_e)$ & $\langle K\rangle$ &
$\langle N_{\mathrm{eff}}\rangle$ & $\langle x_{\max}\rangle$ & $\langle H\rangle$\\
\midrule
1,0 & \direct  & 6,658 & 1,182 & 2,662 & 0,464 & 0,929\\
1,0 & \delayed & 4,651 & 2,121 & 2,367 & 0,530 & 0,842\\
1,3 & \direct  & 7,332 & 1,295 & 2,654 & 0,466 & 0,926\\
1,3 & \delayed & 5,222 & 2,276 & 2,355 & 0,535 & 0,838\\
1,6 & \direct  & 7,705 & 1,417 & 2,633 & 0,471 & 0,920\\
1,6 & \delayed & 5,510 & 2,337 & 2,367 & 0,536 & 0,843\\
1,7 & \direct  & 7,543 & 1,398 & 2,628 & 0,473 & 0,919\\
1,7 & \delayed & 5,384 & 2,383 & 2,374 & 0,533 & 0,846\\
\bottomrule
\end{tabular}
\end{table}

Różnica $\langle K\rangle_{\mathrm{delayed}}-\langle K\rangle_{\mathrm{direct}}$ wynosi około $0{,}92$--$0{,}99$ dla wszystkich natężeń. $K$ zachowuje skalę pędu, natomiast $N_{\mathrm{eff}}$, $x_{\max}$ i $H$ zależą od proporcji. Ich zgodna zmiana pokazuje, że wynik nie jest wyłącznie skutkiem różnicy całkowitej skali.

\section{Sektory znakowe}

\begin{table}[htbp]
\centering
\caption{Udział zdarzeń, w których trzy pędy mają jednakowy znak.}
\label{tab:sign-all}
\begin{tabular}{clrrr}
\toprule
$I$ & kanał & $+++$ [\%] & $---$ [\%] & razem [\%]\\
\midrule
1,0 & \direct  & 48,29 & 49,07 & 97,36\\
1,0 & \delayed & 21,89 & 24,65 & 46,54\\
1,3 & \direct  & 46,77 & 48,52 & 95,29\\
1,3 & \delayed & 24,76 & 23,47 & 48,23\\
1,6 & \direct  & 45,18 & 45,65 & 90,83\\
1,6 & \delayed & 23,81 & 24,44 & 48,25\\
1,7 & \direct  & 47,43 & 43,07 & 90,50\\
1,7 & \delayed & 23,05 & 22,67 & 45,72\\
\bottomrule
\end{tabular}
\end{table}

Dla \direct{} udział sektorów zgodnych przekracza $90\%$ przy każdym natężeniu. Dla \delayed{} pozostaje blisko $46$--$48\%$. Informacja ta zniknęłaby po pozostawieniu wyłącznie modułów pędu.

\section{Korelacje parami}

\begin{table}[htbp]
\centering
\caption{Korelacje Pearsona pomiędzy pędami podłużnymi.}
\label{tab:corr-all}
\small
\begin{tabular}{clrrr}
\toprule
$I$ & kanał & $r_{23}$ & $r_{24}$ & $r_{34}$\\
\midrule
1,0 & \direct  & 0,867 & 0,844 & 0,833\\
1,0 & \delayed & 0,263 & 0,428 & 0,382\\
1,3 & \direct  & 0,871 & 0,830 & 0,834\\
1,3 & \delayed & 0,268 & 0,435 & 0,459\\
1,6 & \direct  & 0,872 & 0,813 & 0,813\\
1,6 & \delayed & 0,269 & 0,470 & 0,465\\
1,7 & \direct  & 0,868 & 0,813 & 0,814\\
1,7 & \delayed & 0,214 & 0,456 & 0,437\\
\bottomrule
\end{tabular}
\end{table}

W \direct{} wszystkie korelacje są silnie dodatnie i stabilne względem natężenia. W \delayed{} korelacja elektronów 2 i 3 jest słabsza, a pozostałe dwie umiarkowane. Wynik jest zgodny z mniejszym $K$, większym $N_{\mathrm{eff}}$ i dominacją sektorów zgodnych w \direct{}, lecz współczynnik Pearsona pozostaje tylko miarą pomocniczą.


\section{Porównanie kanałów direct i delayed dla $I=1{,}7\,\mathrm{PW/cm^2}$}

Aktualne uruchomienie notebooka obejmuje 1695 zdarzeń direct i 1601
zdarzeń delayed. Tabele 5 i 6 zestawiają główne wartości liczbowe.
Histogramy używają surowych liczebności, dlatego przy porównaniach
kształtu warto pamiętać o niewielkiej różnicy liczby zdarzeń; wartości
średnie i współczynniki korelacji nie są obciążone tym efektem w taki
sam sposób.

\begin{longtable}[]{@{}
  >{\raggedright\arraybackslash}p{(\columnwidth - 16\tabcolsep) * \real{0.1111}}
  >{\raggedright\arraybackslash}p{(\columnwidth - 16\tabcolsep) * \real{0.1111}}
  >{\raggedright\arraybackslash}p{(\columnwidth - 16\tabcolsep) * \real{0.1111}}
  >{\raggedright\arraybackslash}p{(\columnwidth - 16\tabcolsep) * \real{0.1111}}
  >{\raggedright\arraybackslash}p{(\columnwidth - 16\tabcolsep) * \real{0.1111}}
  >{\raggedright\arraybackslash}p{(\columnwidth - 16\tabcolsep) * \real{0.1111}}
  >{\raggedright\arraybackslash}p{(\columnwidth - 16\tabcolsep) * \real{0.1111}}
  >{\raggedright\arraybackslash}p{(\columnwidth - 16\tabcolsep) * \real{0.1111}}
  >{\raggedright\arraybackslash}p{(\columnwidth - 16\tabcolsep) * \real{0.1111}}@{}}
\caption{Wielkości kolektywne, względne i korelacje parami.}\label{tab:5}\\
\toprule\noalign{}
\begin{minipage}[b]{\linewidth}\raggedright
\textbf{Kanał}
\end{minipage} & \begin{minipage}[b]{\linewidth}\raggedright
\textbf{n}
\end{minipage} & \begin{minipage}[b]{\linewidth}\raggedright
\textbf{\ensuremath{\langle \Pe{}\rangle}}
\end{minipage} & \begin{minipage}[b]{\linewidth}\raggedright
\textbf{\ensuremath{\sigma}(\Pe{})}
\end{minipage} & \begin{minipage}[b]{\linewidth}\raggedright
\textbf{\ensuremath{\langle K\rangle}}
\end{minipage} & \begin{minipage}[b]{\linewidth}\raggedright
\textbf{\ensuremath{\sigma}(K)}
\end{minipage} & \begin{minipage}[b]{\linewidth}\raggedright
\textbf{\ensuremath{r_{23}}}
\end{minipage} & \begin{minipage}[b]{\linewidth}\raggedright
\textbf{\ensuremath{r_{24}}}
\end{minipage} & \begin{minipage}[b]{\linewidth}\raggedright
\textbf{\ensuremath{r_{34}}}
\end{minipage} \\
\midrule\noalign{}
\endhead
\bottomrule\noalign{}
\endlastfoot
direct & 1695 & 0,320 & 7,543 & 1,398 & 0,665 & 0,868 & 0,813 & 0,814 \\
delayed & 1601 & 0,147 & 5,384 & 2,383 & 1,139 & 0,214 & 0,456 &
0,437 \\
\end{longtable}

\begin{longtable}[]{@{}
  >{\raggedright\arraybackslash}p{(\columnwidth - 14\tabcolsep) * \real{0.1250}}
  >{\raggedright\arraybackslash}p{(\columnwidth - 14\tabcolsep) * \real{0.1250}}
  >{\raggedright\arraybackslash}p{(\columnwidth - 14\tabcolsep) * \real{0.1250}}
  >{\raggedright\arraybackslash}p{(\columnwidth - 14\tabcolsep) * \real{0.1250}}
  >{\raggedright\arraybackslash}p{(\columnwidth - 14\tabcolsep) * \real{0.1250}}
  >{\raggedright\arraybackslash}p{(\columnwidth - 14\tabcolsep) * \real{0.1250}}
  >{\raggedright\arraybackslash}p{(\columnwidth - 14\tabcolsep) * \real{0.1250}}
  >{\raggedright\arraybackslash}p{(\columnwidth - 14\tabcolsep) * \real{0.1250}}@{}}
\caption{Miary równomierności podziału i dominujące sektory znaków.}\label{tab:6}\\
\toprule\noalign{}
\begin{minipage}[b]{\linewidth}\raggedright
\textbf{Kanał}
\end{minipage} & \begin{minipage}[b]{\linewidth}\raggedright
\textbf{\ensuremath{\langle \Neff{}\rangle}}
\end{minipage} & \begin{minipage}[b]{\linewidth}\raggedright
\textbf{med(\Neff{})}
\end{minipage} & \begin{minipage}[b]{\linewidth}\raggedright
\textbf{\ensuremath{\langle \xmax{}\rangle}}
\end{minipage} & \begin{minipage}[b]{\linewidth}\raggedright
\textbf{med(\xmax{})}
\end{minipage} & \begin{minipage}[b]{\linewidth}\raggedright
\textbf{\ensuremath{\langle H\rangle}}
\end{minipage} & \begin{minipage}[b]{\linewidth}\raggedright
\textbf{dominujący sektor}
\end{minipage} & \begin{minipage}[b]{\linewidth}\raggedright
\textbf{udział}
\end{minipage} \\
\midrule\noalign{}
\endhead
\bottomrule\noalign{}
\endlastfoot
direct & 2,628 & 2,726 & 0,473 & 0,456 & 0,919 & \ensuremath{+++} & 47,4\% \\
delayed & 2,374 & 2,386 & 0,533 & 0,518 & 0,846 & \ensuremath{+++} & 23,0\% \\
\end{longtable}

Średni całkowity pęd \Pe{} jest w obu zbiorach bliski zeru w porównaniu z
jego szerokością. Nie oznacza to braku silnego ruchu kolektywnego:
dodatnie i ujemne gałęzie rozkładu częściowo znoszą się w średniej.
Odchylenie standardowe \Pe{} jest większe w kanale direct, co odpowiada
szerszemu rozdzieleniu dwóch gałęzi o przeciwnych znakach.

Najbardziej jednoznaczna różnica dotyczy nierównowagi względnej. Średnia
wartość K jest w kanale delayed większa o 0,985 jednostki pędu.
Jednocześnie średnie \Neff{} jest mniejsze o 0,254, a średnie \xmax{}
większe o 0,060. Wszystkie trzy miary prowadzą do tego samego wniosku: w
zdarzeniach delayed pęd jest przeciętnie dzielony mniej równomiernie, a
pojedynczy elektron częściej przejmuje większą część bezwzględnego pędu.

Znormalizowana entropia udziałów wynosi 0,919 dla direct i 0,846 dla
delayed. Różnica jest zgodna z \Neff{} i \xmax{}, lecz wykorzystuje pełny
wektor udziałów, a nie tylko jego największy element. W kanale direct
blisko połowa zdarzeń należy do sektora \ensuremath{+++}; z wykresu sektorów wynika
również duży udział sektora \ensuremath{---}. W kanale delayed żaden pojedynczy
sektor nie dominuje równie silnie, a sektory mieszane są znacznie
częstsze.

Wycentrowany moment trzeciego rzędu wynosi \ensuremath{-}0,997 dla direct i 0,123 dla
delayed. Różnica sygnalizuje odmienną strukturę wyższych korelacji,
jednak bez normalizacji, błędów statystycznych i analizy zależności od
natężenia nie należy przypisywać jej samodzielnej interpretacji
mechanistycznej.

\chapter{Dyskusja}
\section{Interpretacja fizyczna}

Zestaw obserwabli tworzy spójny obraz kanału direct. Prawie równe
współczynniki PC1, silne dodatnie korelacje parami, dominacja sektorów o
jednakowych znakach, małe K oraz wartości \Neff{} bliskie 3 wskazują na
ruch o charakterze kolektywnym. Taki obraz jest zgodny z mechanizmem, w
którym powracający elektron przekazuje energię pozostałym elektronom w
krótkim przedziale czasu, a wszystkie trzy elektrony opuszczają obszar
jonu w zbliżonej fazie pola. Wspólna faza przyspieszania prowadzi
następnie do zgodnych znaków i podobnych wartości pędu końcowego.

Kanał delayed wykazuje słabsze korelacje, większe K, mniejsze \Neff{},
większe \xmax{} oraz znacznie więcej sektorów mieszanych. Jest to zgodne
z emisją rozdzieloną czasowo: co najmniej jeden elektron pozostaje w
pobliżu jonu dłużej i jest przyspieszany w innej fazie pola niż
pozostałe. W rezultacie końcowe pędy są mniej zgodne i częściej różnią
się zarówno wartością, jak i znakiem.

Powyższa interpretacja ma charakter zgodności z mechanizmem, a nie jego
dowodu. Końcowy rozkład pędów jest projekcją całej dynamiki na stan
asymptotyczny. Rozróżnienie sekwencji czasowej wymaga informacji
trajektoryjnej, czasów jonizacji lub dodatkowych obserwabli. Analiza
pędów dostarcza natomiast zwartego i ilościowego testu, czy dane
sklasyfikowane jako direct i delayed rzeczywiście mają oczekiwaną
strukturę korelacji.

\section{Zalety i ograniczenia metody}

\subsection{Zalety}

\begin{itemize}
\item
  Stały układ współrzędnych. Baza Jacobiego--Helmerta nie zależy od
  konkretnej macierzy kowariancji, dzięki czemu można bezpośrednio
  porównywać kanały i natężenia.
\item
  Pełna odwracalność dla trzech podłużnych pędów. Transformacja
  (\ensuremath{p_2},\ensuremath{p_3},\ensuremath{p_4}) \ensuremath{\leftrightarrow} (Q,\ensuremath{\xi_1},\ensuremath{\xi_2}) nie usuwa informacji.
\item
  Rozdzielenie ruchu kolektywnego i względnego. \Pe{} lub Q opisuje ruch
  wspólny, natomiast \ensuremath{\xi_1}, \ensuremath{\xi_2} i K opisują wewnętrzny podział.
\item
  Połączenie z reprezentacją simpleksową. \Neff{}, \xmax{}, H, sektory
  znaków i ilr zachowują najważniejsze informacje simpleksu w zestawie
  prostych wykresów.
\item
  Dostępność miar permutacyjnie niezmienniczych. K, \Neff{}, \xmax{} i H
  ograniczają zależność wniosków od umownej numeracji identycznych
  elektronów.
\item
  Naturalne rozszerzenie na większą liczbę cząstek. Macierz Helmerta
  można zbudować dla dowolnego N, otrzymując jeden tryb kolektywny i \ensuremath{N-1}
  ortogonalnych trybów względnych.
\end{itemize}

\subsection{Ograniczenia i warunki ostrożnej interpretacji}

\begin{itemize}
\item
  Dla N\textgreater3 pełna przestrzeń względna ma wymiar \ensuremath{N-1}. Sam skalar
  K jest wtedy tylko promieniem i nie zachowuje informacji o kierunku
  nierównowagi. Skalowalność dotyczy konstrukcji współrzędnych i zestawu
  projekcji, a nie magicznej redukcji całej fizyki do dwóch osi.
\item
  Wyniki PCA zależą od centrowania i standaryzacji. Współczynniki z
  różnych procedur przygotowania danych nie muszą być bezpośrednio
  porównywalne.
\item
  Niskowariancyjne PC8 i PC9 mogą być bardziej podatne na szum
  numeryczny i selekcję próbki. Ich znaczenie wynika tutaj z geometrii
  podobszaru z, nie z dużego udziału wyjaśnionej wariancji.
\item
  Współrzędne \ensuremath{\xi_1} i \ensuremath{\xi_2} wykorzystują historię etykiet elektronów. Przy
  porównaniu z eksperymentem, w którym nie identyfikuje się elektronu
  tunelującego, należy stosować symetryzację lub miary niezmiennicze na
  permutacje.
\item
  Histogramy 2D przedstawione w notebooku są histogramami liczebności.
  Do precyzyjnego porównania zbiorów o różnej liczbie zdarzeń należy
  stosować identyczne zakresy, te same krawędzie binów i normalizację do
  gęstości prawdopodobieństwa.
\item
  Średnie i korelacje powinny być uzupełnione przedziałami ufności, na
  przykład wyznaczonymi metodą bootstrap. Jest to szczególnie ważne dla
  momentów wyższego rzędu i mało licznych sektorów znaków.
\item
  Wartość K zależy od skali pędów. Przy porównywaniu różnych natężeń
  pola należy zestawiać ją z miarami bezwymiarowymi lub rozważyć
  znormalizowany odpowiednik K/A.
\end{itemize}

\chapter{Podsumowanie}

Analiza głównych składowych ujawniła, że podłużna dynamika trzech
elektronów ma naturalną strukturę składającą się z jednego trybu
kolektywnego i dwóch trybów względnych. Najważniejszym krokiem
metodologicznym było odejście od bezpośredniego używania osi PCA na
rzecz stałej, ortonormalnej bazy Jacobiego--Helmerta. Dzięki temu
zachowano fizyczną interpretowalność, pełną odwracalność transformacji
oraz możliwość porównywania niezależnych zbiorów w identycznym układzie
współrzędnych.

Współrzędne Q, \ensuremath{\xi_1} i \ensuremath{\xi_2} rozdzielają sumę pędów od ich wzajemnych różnic.
Moduł K stanowi prostą, permutacyjnie niezmienniczną miarę nierównowagi.
Obserwable simpleksowe \Neff{}, \xmax{} i H opisują równomierność
znormalizowanego podziału, sektory znaków zachowują kierunek emisji, a
transformacja ilr umożliwia czytelną dwuwymiarową prezentację udziałów
bezwzględnych. Żadna z tych wielkości pojedynczo nie zastępuje pełnego
rozkładu; ich siła wynika z komplementarności.

Dla natężenia \ensuremath{1{,}7\,\mathrm{PW/cm^2}} kanał direct wykazuje silne dodatnie korelacje
wszystkich par elektronów, częstą emisję w tym samym kierunku, mniejszą
nierównowagę K oraz bardziej równomierny podział bezwzględnego pędu.
Kanał delayed jest szerszy w przestrzeni względnej, ma słabsze
korelacje, większy udział sektorów mieszanych i większą dominację
pojedynczego elektronu. Zestaw tych wyników odróżnia oba kanały w sposób
znacznie bardziej jednoznaczny niż pojedynczy wykres simpleksowy lub
pojedyncza projekcja PCA.

Metoda stanowi więc logiczne rozwinięcie reprezentacji simpleksowej, a
nie jej prosty zamiennik. Dla trzech elektronów pozwala zachować całą
informację podłużną w trzech współrzędnych oraz rozłożyć jej
interpretację na kilka przejrzystych wykresów. Dla większej liczby
elektronów dostarcza systematycznej bazy wielu ciał i zestawu skalarnych
podsumowań, przy czym pełna analiza nadal wymaga odpowiedniej liczby
współrzędnych względnych lub świadomie dobranych projekcji.

\chapter*{Materiały robocze}
\addcontentsline{toc}{chapter}{Materiały robocze}

\begin{itemize}
\item
  05\_PCA.ipynb --- analiza głównych składowych, współczynniki kierunków
  i pierwsza konstrukcja trzech kombinacji podłużnych pędów.
\item
  07\_analiza\_3D.ipynb --- implementacja współrzędnych
  Jacobiego--Helmerta, obserwabli simpleksowych, testów rekonstrukcji
  oraz porównania direct--delayed.
\item
  praca\_licencjacka\_maciej.pdf --- materiał wykorzystany wyłącznie
  jako odniesienie do stylu redakcyjnego i sposobu prowadzenia narracji
  naukowej.
\end{itemize}

\end{document}
